\chapter{Marge}\label{ch:m1:margin}

Un fonds boucle sa journée du vendredi avec quatre cents contrats à terme sur indice en position longue, entièrement couverts par des marges.
Le lundi, il détient les mêmes quatre cents contrats et doit à la
chambre de compensation des liquidités deux fois~: une fois pour ce que les contrats ont perdu, et une nouvelle fois parce que
la chambre a décidé, en regardant les mêmes cours, que chaque contrat
est désormais plus risqué à détenir. En mars 2020, la \omterm{def:m1:clearing-and-settlement:margin}{marge initiale}
exigée par les \omterm{def:m1:clearing-and-settlement:ccp}{contreparties centrales} du monde entier a augmenté d'environ 300~milliards de dollars,
et, en une seule journée, le 9 mars, elles ont appelé 140~milliards de dollars de \omterm{def:m1:clearing-and-settlement:margin}{marge de variation}.
Aucune de ces sommes n'était une perte pour le système pris dans son ensemble~; toutes devaient être
trouvées en liquidités dans les heures qui suivaient, par quelqu'un. Le \Cref{ch:m1:clearing-and-settlement}
a expliqué à quoi sert la marge. Ce chapitre la calcule comme le fait la
chambre de compensation, et examine comment ce calcul se comporte le jour où il compte le plus.

\section{Marge de variation}

La \omterm{def:m1:clearing-and-settlement:margin}{marge de variation} (\cref{def:m1:clearing-and-settlement:margin}) est
une arithmétique simple~: la position multipliée par la variation du cours de règlement et par le
multiplicateur, payée en liquidités chaque jour, dans la devise du contrat. Trois
propriétés la rendent plus délicate qu'il n'y paraît. Elle est \emph{asymétrique en
liquidité}~: un portefeuille couvert dont l'autre jambe n'est pas auprès de la même
chambre de compensation paie des liquidités sur la jambe perdante et n'en reçoit aucune sur la gagnante. Elle est
\emph{rapide}~: un \omterm{def:m1:financing:margincall}{appel de marge} lancé le matin est exigible le matin même, et sur les
marchés sous tension, les chambres de compensation en relancent dans la journée. Et elle est
\emph{impitoyable}~: un membre qui ne paie pas est en défaut.

\begin{definition}[Dépôt de garantie et marge de maintenance]\label{def:m1:margin:bond}
Un \emph{dépôt de garantie}\index{dépôt de garantie} est le nom que l'industrie des contrats à
terme donne à la \omterm{def:m1:clearing-and-settlement:margin}{marge initiale}~: un collatéral déposé pour garantir l'exécution, et non
un acompte. Pour un compte client, le courtier fixe un niveau \emph{initial},
requis pour ouvrir une position, et un seuil inférieur, la \emph{marge
de maintenance}\index{marge de maintenance}~: lorsque la valeur du compte passe sous le
seuil de maintenance, le client est appelé à la rétablir au niveau
initial. Pour les membres compensateurs, les deux niveaux coïncident.
\end{definition}

\begin{definition}[Appel de marge intrajournalier]\label{def:m1:margin:intraday}
Un \emph{appel de marge intrajournalier}\index{appel de marge intrajournalier} est un appel de
\omterm{def:m1:clearing-and-settlement:margin}{marge de variation} ou \omterm{def:m1:clearing-and-settlement:margin}{marge initiale}, lancé par une chambre de compensation entre ses
cycles quotidiens programmés, sur la base des cours et des positions du moment, exigible
dans l'heure environ ou le délai que fixent ses règles.
\end{definition}

\section{Marge initiale par scénarios}

\begin{definition}[Intervalle de balayage et tableau de risque]\label{def:m1:margin:scan}
Dans un système de marge fondé sur des scénarios, la chambre de compensation fixe, pour chaque produit,
un \emph{intervalle de balayage des prix} et un \emph{intervalle de balayage de la volatilité}~: les plus
grandes variations des prix et de la volatilité implicite qu'elle souhaite couvrir sur son
horizon, soit l'\emph{intervalle de balayage}\index{intervalle de balayage}. Le \emph{tableau de
risque} d'un contrat est son gain ou sa perte dans chacun d'une liste fixe de
scénarios construits à partir de ces intervalles~; le \emph{risque de balayage} d'un portefeuille est sa
plus forte perte totale sur la liste.
\end{definition}

Le système qui a assuré la marge des contrats à terme pendant trois décennies, SPAN, utilise
seize scénarios. Deux ne font varier que la volatilité, à la hausse et à la baisse. Douze combinent une
variation de prix d'un tiers, de deux tiers ou de la totalité de l'\omterm{def:m1:margin:scan}{intervalle de balayage} des prix, à la hausse
ou à la baisse, la volatilité pouvant monter ou baisser. Les deux dernières sont des \emph{mouvements extrêmes}, un
multiple de l'\omterm{def:m1:margin:scan}{intervalle de balayage}, dont seule une fraction de la perte est retenue~:
elles s'appliquent aux options vendues très hors de la monnaie, qui ne perdent rien dans les
quatorze premiers. Pour un contrat à terme, l'\omterm{def:m1:margin:scan}{intervalle de balayage} des prix est la \omterm{def:m1:margin:bond}{marge de
maintenance} elle-même.

\begin{example}[Le tableau de risque d'un contrat à terme]\label{ex:m1:margin:array}
Un examen du système par le régulateur publie le tableau de risque du grand contrat à terme S\&P 500
à terme du 12 avril 2001, lorsque sa \omterm{def:m1:margin:bond}{marge de maintenance} valait 17\,250~\$~: zéro dans
les deux scénarios de volatilité~; pertes de $\mp\$5\,750$, $\mp\$11\,500$ et
$\mp\$17\,250$ à un, deux et trois tiers~; et dans les scénarios extrêmes, un
mouvement égal à trois fois l'intervalle, avec 30~\% de la perte retenue, $\pm\$15\,525$. Le
\omterm{def:m1:margin:scan}{risque de balayage} d'un contrat à terme acheté est le mouvement baissier complet, 17\,250~\$~; le
scénario extrême, à 15\,525~\$, n'est pas contraignant. Il l'est pour le vendeur d'options.
\end{example}

\begin{method}[Une marge par scénarios]\label{met:m1:margin:span}
Pour chaque produit~:
\begin{enumerate}
\item \textbf{\omterm{def:m1:margin:scan}{Risque de balayage}}~: sommer la position multipliée par le tableau de risque sur tous les contrats,
scénario par scénario~; retenir la perte la plus forte.
\item \textbf{Charge inter-mois}~: les contrats à terme d'échéances différentes se compensent
intégralement dans le balayage~; ajouter une charge par spread pour le risque qu'ils ne
évoluent pas de concert.
\item \textbf{Minimum pour les options vendues}~: un plancher par option vendue, comparé à
la somme des deux premiers termes.
\end{enumerate}
Entre produits~: retrancher un \emph{crédit de spread
inter-produits}\index{crédit de spread inter-produits} pour les couples dont les prix sont
corrélés et dont les positions se compensent. Additionner les résultats.
\end{method}

\begin{omfigure}
\begin{tikzpicture}
\begin{axis}[width=11.5cm,height=5cm,ybar,bar width=8pt,
  xlabel={scénario},ylabel={perte du portefeuille (milliers de \$)},
  tick label style={font=\small},label style={font=\small},
  xmin=0.3,xmax=16.7,xtick={1,2,3,4,5,6,7,8,9,10,11,12,13,14,15,16},ymin=-50,ymax=140,ymajorgrids,grid style={gray!25},table/col sep=comma]
\addplot[fill=omDef!60,draw=omDef] table[x=scenario,y=loss_k]{figdata/markets-1/20-margin/scenarios.csv};
\node[font=\footnotesize,anchor=south east] at (axis cs:16.6,112) {risque de balayage~: 107.8};
\end{axis}
\end{tikzpicture}

\omcaption{Seize scénarios pour un portefeuille long de 2 contrats à terme à 6\,000, vendeur de
10 puts de strike 5\,600 et vendeur de 10 calls à 6\,400 (trois mois, 18~\%
volatilité, \omterm{def:m1:margin:scan}{intervalle de balayage} des prix de 345 points, \omterm{def:m1:margin:scan}{intervalle de balayage} de la volatilité de 4 points, mouvements extrêmes
de trois intervalles, retenus à 30~\%). Les scénarios impairs relèvent la volatilité, les pairs
la font baisser. La marge est fixée par le scénario 16, la chute extrême. Données~: la
construction du chapitre.}\label{fig:m1:margin:scenarios}
\end{omfigure}

\begin{omfigure}
\begin{tikzpicture}
\begin{axis}[width=11cm,height=5.2cm,
  xlabel={variation du contrat à terme (points)},ylabel={P\&L (milliers de \$)},
  tick label style={font=\small},label style={font=\small},
  xmin=-1100,xmax=1100,xtick={-1000,-500,0,500,1000},ymin=-420,ymax=80,grid=major,grid style={gray!25},
  legend columns=3,legend style={font=\footnotesize,at={(0.5,-0.3)},anchor=north,draw=none,fill=none,/tikz/every even column/.append style={column sep=8pt}},
  table/col sep=comma]
\addplot[omThm,very thick] table[x=move,y=vol_up_k]{figdata/markets-1/20-margin/profile.csv};
\addplot[omDef,very thick,dashed] table[x=move,y=vol_down_k]{figdata/markets-1/20-margin/profile.csv};
\addplot[only marks,mark=*,mark size=1.8pt,black] table[x=move,y=pnl_k]{figdata/markets-1/20-margin/points.csv};
\draw[thin] (axis cs:-345,-420) -- (axis cs:-345,80);
\draw[thin] (axis cs:345,-420) -- (axis cs:345,80);
\legend{volatilité $+4$ points,volatilité $-4$ points,scénarios 1 à 14}
\end{axis}
\end{tikzpicture}

\omcaption{Ce que voit le balayage. Le même portefeuille réévalué sur une large plage
de prix~; les quatorze scénarios standard sont les points, tous situés à l'intérieur de
l'\omterm{def:m1:margin:scan}{intervalle de balayage} des prix (traits fins). Les pertes qui comptent pour ce portefeuille
se situent hors de celui-ci, et c'est à cela que servent les scénarios 15 et 16.
}\label{fig:m1:margin:profile}
\end{omfigure}

\begin{definition}[Marge en portefeuille]\label{def:m1:margin:portfolio}
La \emph{marge en portefeuille}\index{marge en portefeuille} (portfolio margining) fixe l'exigence sur le
risque net de l'ensemble des positions d'un compte, en tenant compte des
compensations entre elles, par opposition à une somme d'exigences calculées
position par position.
\end{definition}

\begin{omfigure}
\begin{tikzpicture}
\begin{axis}[width=9.5cm,height=4.4cm,ybar,bar width=16pt,
  ylabel={marge (milliers de \$)},
  symbolic x coords={future,short put,short call,portfolio},xtick={future,short put,short call,portfolio},xticklabels={{future},{put vendu},{call vendu},{portefeuille}},
  tick label style={font=\small},label style={font=\small},xticklabel style={font=\small\strut},
  ymin=0,ymax=135,ymajorgrids,grid style={gray!25},enlarge x limits=0.18,
  nodes near coords,nodes near coords style={font=\footnotesize,/pgf/number format/.cd,fixed,fixed zerofill,precision=1},
  table/col sep=comma]
\addplot[fill=omDef!60,draw=omDef] table[x=item,y=margin_k]{figdata/markets-1/20-margin/offsets.csv};
\end{axis}
\end{tikzpicture}

\omcaption{Compensations. Margés séparément, les trois jambes exigent 213\,000~\$~;
ensemble, 108\,000~\$, car le put et le call ne peuvent pas perdre tous deux. Données~:
la construction du chapitre.}\label{fig:m1:margin:offsets}
\end{omfigure}

\section{Marge initiale par la valeur en risque}

Seize scénarios par produit et un tableau de crédits entre produits sont
transparents et reproductibles, ce qui explique que la méthode ait duré. Ils sont aussi
grossiers~: les compensations sont fixées par un comité, et un portefeuille de nombreux produits est
margé comme une somme de paires. Les grandes chambres de compensation sont passées, ou
passent, à la valeur en risque de portefeuille.

\begin{definition}[Période de marge en risque]\label{def:m1:margin:mpor}
La \emph{période de marge en risque}\index{période de marge en risque} (margin period of risk) est le délai
supposé entre le dernier paiement de \omterm{def:m1:clearing-and-settlement:margin}{marge de variation} d'un membre compensateur et le moment
où son portefeuille a été couvert ou liquidé par la chambre de compensation~:
l'horizon du calcul de la \omterm{def:m1:clearing-and-settlement:margin}{marge initiale}. Il est d'un jour ou deux pour les contrats
à terme cotés liquides, et plus long pour les swaps compensés et les produits peu liquides.
\end{definition}

\begin{dated}{2026-09}{Des scénarios à la simulation}\label{dat:m1:margin:models}
Le successeur du SPAN à la chambre de compensation de Chicago repose sur la valeur en
risque historique~: des milliers de scénarios issus d'au moins dix ans de données, avec un
redimensionnement de la volatilité et des corrélations, et une présentation séparée des
composantes de marché, de liquidité et de concentration. La méthode de portefeuille de la chambre
de compensation de Francfort repose sur la simulation historique filtrée, associée à une
valeur en risque en période de tensions, calculée par groupes de produits qui seraient
liquidés ensemble, chacun avec sa propre \omterm{def:m1:margin:mpor}{période de marge en risque}.
\end{dated}

\begin{method}[Simulation historique filtrée]\label{met:m1:margin:fhs}
(i) Estimer la volatilité $\sigma_t$ de chaque jour à partir des rendements passés, par
exemple par une pondération exponentielle. (ii) Standardiser les $n$ derniers rendements~:
$z_s = r_s/\sigma_s$. (iii) Les ramener à aujourd'hui~: $\tilde r_s = z_s\,\sigma_t$.
(iv) Réévaluer le portefeuille sous chaque $\tilde r_s$, prendre le quantile de perte au
niveau de confiance retenu, et le ramener à la \omterm{def:m1:margin:mpor}{période de marge en risque}.
\end{method}

Le filtrage fait réagir la marge en quelques jours à un changement de régime, là où une
fenêtre historique simple ne réagit qu'à mesure que les journées extrêmes s'accumulent. C'est
la bonne propriété pour couvrir la perte de demain, et la mauvaise pour la stabilité du
système.

\section{Procyclicité}

\begin{definition}[Procyclicité]\label{def:m1:margin:procyclical}
Un système de marge est \emph{procyclique}\index{procyclicité} (procyclical) dans la mesure où
ses exigences augmentent sur les marchés sous tension et diminuent sur les marchés calmes,
ajoutant à la demande de liquidité précisément quand celle-ci est rare.
\end{definition}

\begin{omfigure}
\begin{tikzpicture}
\begin{axis}[width=11cm,height=5.2cm,
  xlabel={jours de négociation depuis le début de la tension},ylabel={marge (\% de la position)},
  tick label style={font=\small},label style={font=\small},
  xmin=-100,xmax=299,ymin=0,ymax=10.5,grid=major,grid style={gray!25},
  legend columns=3,legend style={font=\footnotesize,at={(0.5,-0.3)},anchor=north,draw=none,fill=none,/tikz/every even column/.append style={column sep=8pt}},
  table/col sep=comma]
\addplot[omExo,very thick] table[x=day,y=hs]{figdata/markets-1/20-margin/procyclical.csv};
\addplot[omThm,very thick] table[x=day,y=fhs]{figdata/markets-1/20-margin/procyclical.csv};
\addplot[omDef,very thick,dashed] table[x=day,y=floored]{figdata/markets-1/20-margin/procyclical.csv};
\legend{simulation historique,filtrée,filtrée avec un plancher en tension de 25~\%}
\end{axis}
\end{tikzpicture}

\omcaption{Marge sur une position fixe pendant que la volatilité quotidienne passe de
0.8~\% à 3.5~\% pendant 25 jours puis décroît (99~\%, horizon de deux jours,
fenêtre de 250 jours). Sur les dix pires jours, la marge brute augmente de 109~\%, la
marge filtrée de 203~\%, la marge plancherisée de 94~\%. Le plancher se paie
en temps calme~: 4.5~\% de la position contre 3.0~\%. Données~: la simulation du
tutoriel.}\label{fig:m1:margin:procyclical}
\end{omfigure}

La revue conjointe de mars 2020, menée par les organismes de normalisation bancaire, des
infrastructures de marché et des valeurs mobilières, a constaté que la hausse de la \omterm{def:m1:clearing-and-settlement:margin}{marge
initiale} compensée était forte et inégale selon les classes d'actifs et les chambres de
compensation, et a inscrit parmi ses chantiers la transparence de la marge compensée, la
préparation des intervenants de marché face à la liquidité, et la réactivité des modèles de
\omterm{def:m1:clearing-and-settlement:margin}{marge initiale} face aux tensions. Pour une société de trading, la conclusion opérationnelle
n'attend pas la politique~: la liquidité nécessaire pour tenir une position est sa \omterm{def:m1:clearing-and-settlement:margin}{marge de
variation} d'une semaine difficile \emph{plus} l'augmentation de sa \omterm{def:m1:clearing-and-settlement:margin}{marge initiale} pendant
cette même semaine, et les deux doivent être modélisées.

\section{Tutoriel~: une marge de scénarios pour contrats à terme et options}\label{tut:m1:margin:span}

\begin{tutorial}
\textbf{Objectif.} Construire les tableaux de risque des contrats à terme et des options sur contrats à terme, balayer
un portefeuille, ajouter les frais, et comparer une marge de valeur en risque filtrée à une marge non filtrée lors
d'un choc de volatilité. \textbf{État final~:} les quatre figures de données
de ce chapitre.
\begin{tutsteps}
\item \textbf{Seize scénarios} construits à partir de deux intervalles et de deux paramètres de
mouvements extrêmes.
\omcode{code/firm/margin/firm_margin.py}{53}{60}{La liste des scénarios~: variation du prix, variation de la volatilité, fraction de la perte prise en compte.}\label{lst:m1:margin:scenarios}
\item \textbf{Un tableau de risque} par réévaluation complète, avec la formule de Black (1976)
pour les options sur contrats à terme.
\omcode{code/firm/margin/firm_margin.py}{63}{74}{Réévaluation et tableau de risque d'une unité longue.}\label{lst:m1:margin:array}
\item \textbf{Balayage, charge, plancher.}
\omcode{code/firm/margin/firm_margin.py}{86}{102}{La marge des positions d'un produit.}\label{lst:m1:margin:product}
\item \textbf{\omterm{def:m1:margin:procyclical}{Procyclicité}.} Recalibrer la fenêtre sur la volatilité du jour.
\omcode{code/markets-1/20-margin/python/margin_demo.py}{30}{33}{Simulation historique filtrée.}\label{lst:m1:margin:fhs}
\end{tutsteps}
\textbf{Pour aller plus loin.} Remplacer les deux scénarios extrêmes par une grille complète
allant jusqu'à cinq intervalles, et voir quels portefeuilles étaient sous-margés. Puis donner
à la méthode filtrée un plancher sur $\sigma_t$ au lieu d'une majoration de stress,
et comparer les deux outils à \omterm{def:m1:pnl-and-positions:realised}{coût moyen} égal.
\end{tutorial}

\section{Construction~: le moteur de marge}\label{bld:m1:margin:engine}

\begin{build}
\bfield{Rôle} La composante de compensation de la firme miniature
(\cref{ch:m1:clearing-and-settlement}) a besoin d'un chiffre chaque nuit et sur demande~: la \omterm{def:m1:clearing-and-settlement:margin}{marge initiale}
de chaque compte, recalculable par le membre. \bfield{Interface} \texttt{Params(price\_scan, vol\_scan, extreme\_mult,
extreme\_cover, intermonth\_charge, short\_option\_min, multiplier)}~;
\texttt{Position(root, expiry, qty, strike, right)}; \texttt{Market};
\texttt{scenarios(params)}; \texttt{risk\_array(position, market, params)};
\texttt{product\_margin(positions, market, params)} renvoie le \omterm{def:m1:margin:scan}{risque de balayage},
le scénario contraignant, la charge intermois, le minimum sur options vendues
et le total~; \texttt{portfolio\_margin(by\_root, market, params, credits)}.
\bfield{Règles} Données des contrats reprises de
la \cref{bld:m1:futures-contracts-and-exchanges:master}. Réévaluation complète, sans approximation par le delta.
Les contrats à terme d'échéances différentes se compensent dans le balayage et paient la charge
intermois par spread. Le \omterm{met:m1:margin:span}{crédit de spread inter-produits} n'est accordé que lorsque les scénarios contraignants
des deux produits sont des variations de prix
de sens opposés.
\bfield{Tests d'acceptation} \texttt{code/firm/margin/tests/}~: le tableau du contrat à terme
de l'\cref{ex:m1:margin:array}, au dollar près~; un \omterm{def:m1:matching-algorithms-and-implied-spreads:calendar}{calendar spread}~; une position vendeuse sur un put lointain hors de la monnaie,
captée par le scénario extrême, et une autre captée par le minimum~; une option couverte~;
des crédits accordés et refusés.
\bfield{Pour aller plus loin} Un second moteur, la valeur en risque de portefeuille par
simulation historique filtrée, derrière la même interface, et un rapport de l'écart
entre les deux sur chaque compte.
\end{build}

\begin{omsources}
\item US Commodity Futures Trading Commission, \emph{Review of Standard
Portfolio Analysis of Risk (``SPAN'') Margin System}, avril 2001.
\item Basel Committee on Banking Supervision, Committee on Payments and
Market Infrastructures and IOSCO, \emph{Review of margining practices},
rapport final et communiqué de presse, 29 septembre 2022.
\item B.~H.~Cohen and K.~Tracol, ``Market turbulence and soaring margins:
lessons from two recent episodes'', \emph{BIS Quarterly Review}, mars 2023.
\item CME Group, \emph{SPAN 2 methodology}; Eurex Clearing, \emph{Eurex
Clearing Prisma}.
\item F.~Black, ``The pricing of commodity contracts'', \emph{Journal of
Financial Economics} 3 (1976).
\end{omsources}

\section{Exercices}

\begin{exercise}[$\star$]\label{exo:m1:margin:1}
Un compte est long de 25 E-mini. Le cours de règlement passe de 6\,012.50 à
5\,968.25. Calculer la \omterm{def:m1:clearing-and-settlement:margin}{marge de variation}, en indiquant son sens.
\end{exercise}

\begin{exercise}[$\star$]\label{exo:m1:margin:2}
Un courtier exige 23\,100~\$ de \omterm{def:m1:clearing-and-settlement:margin}{marge initiale} et 21\,000~\$ de \omterm{def:m1:margin:bond}{marge de maintenance} par
contrat. Un client détient 5 contrats avec 120\,000~\$ de capitaux propres et perd
3\,200~\$ par contrat. Y a-t-il un \omterm{def:m1:financing:margincall}{appel de marge}, et pour quel montant~?
\end{exercise}

\begin{exercise}[$\star$]\label{exo:m1:margin:3}
Écrire le tableau de risque à seize entrées d'un contrat à terme long, de multiplicateur
50~\$, avec un \omterm{def:m1:margin:scan}{intervalle de balayage} du prix de 345 points et des mouvements extrêmes de trois intervalles
comptés à 30~\%. Quelle entrée constitue sa marge~?
\end{exercise}

\begin{exercise}[$\star\star$]\label{exo:m1:margin:4}
Dans la \cref{fig:m1:margin:scenarios}, pourquoi les scénarios impairs font-ils perdre plus que
leurs voisins pairs~? Pourquoi le scénario 16 est-il pire que le scénario 15 alors que le
strangle est symétrique~? Qu'est-ce qui ferait du scénario 13 le scénario contraignant~?
\end{exercise}

\begin{exercise}[$\star\star$]\label{exo:m1:margin:5}
À partir de la \cref{fig:m1:margin:offsets}, calculer l'économie réalisée en margeant
le portefeuille globalement, en dollars et en pourcentage. Un courtier marge chaque jambe
séparément «~par prudence~». Que coûte cette prudence à un client
qui paie 5~\% par an de financement~?
\end{exercise}

\begin{exercise}[$\star\star$]\label{exo:m1:margin:6}
Une valeur en risque à un jour au niveau de 99~\% vaut 2.1~\% d'une position. Calculer la marge
pour des \omterm{def:m1:margin:mpor}{périodes de marge en risque} de deux et de cinq jours selon la mise à l'échelle en racine carrée.
Donner une raison pour laquelle le chiffre réel à cinq jours pourrait être plus élevé, et une raison
pour laquelle il pourrait être plus bas.
\end{exercise}

\begin{exercise}[$\star\star\star$]\label{exo:m1:margin:7}
\emph{Programmation.} Avec \texttt{peak\_to\_trough} et \texttt{procyclical.csv},
reproduire les trois hausses sur dix jours citées dans
la \cref{fig:m1:margin:procyclical} et les deux moyennes des périodes calmes. Trouver le
poids du plancher de stress pour lequel la hausse sur dix jours vaut 50~\%, et son
coût en périodes calmes.
\end{exercise}

\begin{exercise}[$\star\star\star$]\label{exo:m1:margin:8}
\emph{Trouver la faille.} «~Notre livre est neutre au marché~: long de 200 millions de dollars
d'actions chez le \omterm{def:m1:financing:pb}{prime broker}, short de 200 millions de dollars de contrats à terme sur indices à la
chambre de compensation. Une baisse de 10~\% du marché ne nous coûte rien, donc les \omterm{def:m1:financing:margincall}{appels de marge}
ne sont pas un risque pour nous.~» Et une hausse de 10~\%~?
\end{exercise}

\section{Problème~: un portefeuille fixe dans un mois comme mars 2020}

\begin{problem}[{Problème du week-end --- mêmes contrats, plus de liquidités}]\label{pb:m1:margin:1}
Les chiffres sont inventés, dans les proportions d'un mois violent. Un fonds est
long de 400 contrats à terme sur indices (multiplicateur 50~\$). En début de mois, le
contrat à terme vaut 3\,300 et la \omterm{def:m1:clearing-and-settlement:margin}{marge initiale} de la chambre de compensation est de 4.5~\% du
notionnel. Quatre semaines plus tard, le contrat à terme vaut 2\,500 et la marge vaut 9~\% du
notionnel. Le fonds conserve 12 millions de dollars de liquidités et détient le reste de ses
actifs en obligations d'entreprises.

\medskip\noindent\textbf{Partie I --- Deux appels.}
\begin{enumerate}
\item Donner le notionnel et la \omterm{def:m1:clearing-and-settlement:margin}{marge initiale} en début de mois.
\item Donner la \omterm{def:m1:clearing-and-settlement:margin}{marge de variation} versée au cours du mois.
\item Donner la \omterm{def:m1:clearing-and-settlement:margin}{marge initiale} en fin de mois, et sa variation.
\item Donner le total des liquidités versées, en dollars et en pourcentage du
notionnel initial.
\item Comparer avec les liquidités du fonds. Que doit-il faire~?
\end{enumerate}

\medskip\noindent\textbf{Partie II --- Le pire jour.}
Un jour, le contrat à terme chute de 9.5~\% à partir de 2\,750.
\begin{enumerate}[resume]
\item Donner la \omterm{def:m1:clearing-and-settlement:margin}{marge de variation} de ce jour.
\item La chambre de compensation émet un \omterm{def:m1:margin:intraday}{appel de marge intrajournalier} à midi, alors que le contrat à terme
a baissé de 6~\%. Donner le montant et la conséquence du délai.
\item Le taux de marge est relevé le soir même de 6~\% à 7.5~\% du notionnel. Donner la \omterm{def:m1:clearing-and-settlement:margin}{marge initiale}
supplémentaire au nouveau cours.
\item Donner le total des liquidités appelées en vingt-quatre heures.
\end{enumerate}

\medskip\noindent\textbf{Partie III --- Modèles.}
\begin{enumerate}[resume]
\item À l'aide de la \cref{fig:m1:margin:procyclical}, quel modèle aurait appelé le plus de marge supplémentaire
au cours des dix premiers jours, et lequel le moins~?
\item Quel modèle protège le mieux la chambre de compensation au quinzième jour~?
\item Expliquer pourquoi le modèle avec plancher est impopulaire les années calmes.
\item Un fonds peut choisir entre deux courtiers~: l'un répercute la marge de la
chambre de compensation, l'autre la facture 1.5 fois en permanence. Lequel est le plus sûr
pour le fonds en période de crise~?
\item La chambre de compensation double les marges. Montrer, en une phrase de calcul,
pourquoi le système dans son ensemble a besoin de plus de collatéral alors que personne
n'est devenu plus pauvre.
\end{enumerate}

\medskip\noindent\textbf{Partie IV --- Jugement.}
\begin{enumerate}[resume]
\item Les contrats à terme du fonds couvraient une position courte en actions détenues
ailleurs. Cela aide-t-il face aux appels~?
\item Que le fonds aurait-il dû mesurer avant le début du mois~?
\item Pourquoi le fonds ne peut-il pas simplement déposer ses obligations d'entreprises~?
\item Si de nombreux fonds vendent des obligations pour faire face aux appels, que se passe-t-il ensuite~?
\item Énoncer le \emph{résultat nommé}~: la \omterm{def:m1:clearing-and-settlement:margin}{marge initiale} supplémentaire appelée sur
le portefeuille inchangé au cours du mois.
\item En une phrase~: quelle est la différence entre être solvable et
être capable de payer un \omterm{def:m1:financing:margincall}{appel de marge}~?
\end{enumerate}
\end{problem}

\section{Questions d'entretien}

\begin{interviewq}[$\star$ \iqroles{trader, chercheur, développeur}]\label{iq:m1:margin:1}
Quelle est la différence entre la \omterm{def:m1:clearing-and-settlement:margin}{marge initiale} et la \omterm{def:m1:clearing-and-settlement:margin}{marge de variation}~?
\end{interviewq}

\begin{interviewq}[$\star$ \iqroles{trader, banque}]\label{iq:m1:margin:2}
Décrivez comment un système de marge fondé sur des scénarios, comme SPAN, calcule
l'exigence pour un portefeuille de contrats à terme et d'options.
\end{interviewq}

\begin{interviewq}[$\star\star$ \iqroles{chercheur, banque}]\label{iq:m1:margin:3}
Qu'est-ce que la \omterm{def:m1:margin:procyclical}{procyclicité} de la marge, et quels outils la réduisent~? Quel est leur
coût~?
\end{interviewq}

\begin{interviewq}[$\star\star$ \iqroles{trader, chercheur}]\label{iq:m1:margin:4}
Vous êtes vendeur d'options très hors de la monnaie. Quelle partie du calcul de la marge
vous vise-t-elle, et pourquoi existe-t-elle~?
\end{interviewq}

\begin{interviewq}[$\star\star$ \iqroles{chercheur, ingénieur ML}]\label{iq:m1:margin:5}
Comment intégreriez-vous la marge dans un backtest d'une stratégie sur contrats à terme à effet de levier~?
\end{interviewq}

\begin{interviewq}[$\star\star\star$ \iqroles{trader, banque, chercheur}]\label{iq:m1:margin:6}
Votre société est couverte sur deux chambres de compensation. Décrivez son risque de liquidité
et comment vous dimensionneriez le coussin de trésorerie.
\end{interviewq}
